\section{A unified map of control representations}
\label{sec:unified-coordinate-map}

Field-oriented control names the alignment principle; current-vector control
usually closes \(dq\) current loops inside it. Universal field orientation
seeks a common flux-aligned interface across machine types
\cite{profumo1992ufo,casadei2013unified}. DTC selects voltage actions from
torque/flux errors, while DFVC regulates continuous flux-oriented quantities.
Active flux changes the torque-producing flux variable; feedback
linearization and deadbeat control change how voltages are computed.
Task/null-space control organizes priorities in state space. Natural gradients
choose locally efficient directions, and hyperbolic/projective coordinates
parameterize torque-factor allocation. These are intersecting design choices,
not mutually exclusive product names.

\begin{figure}[tb]
\centering
\begin{tikzpicture}[font=\scriptsize,>=Latex,node distance=6mm and 7mm,
 box/.style={draw,rounded corners,align=center,minimum height=7mm,fill=MidnightBlue!6}]
\node[box] (foc) {FOC / CVC\\rotor \(dq\)};
\node[box,right=of foc] (ufo) {UFO / DFVC\\flux frame};
\node[box,right=of ufo] (dtc) {DTC / deadbeat\\direct voltage action};
\node[box,below=of foc] (af) {active flux\\torque flux factor};
\node[box,right=of af] (task) {task + null space\\priorities};
\node[box,right=of task] (iofl) {feedback linearization\\virtual inputs};
\node[box,below=of af] (hyp) {hyperbolic/projective\\allocation chart};
\node[box,right=of hyp] (grad) {natural gradient\\local metric frame};
\node[box,right=of grad] (ajd) {AJD\\fixed compromise};
\draw[<->] (foc)--(ufo);
\draw[<->] (ufo)--(dtc);
\draw[->] (af)--(task);
\draw[->] (task)--(iofl);
\draw[->] (hyp)--(task);
\draw[->] (grad)--(task);
\draw[->] (ajd)--(task);
\draw[->] (ufo)--(af);
\draw[->] (dtc)--(iofl);
\end{tikzpicture}
\caption{Geometric map of control families. Spatial frames, torque variables,
task priorities, and voltage laws are composable layers rather than mutually
exclusive alternatives.}
\label{fig:control-concept-map}
\end{figure}


\begin{landscape}
\begin{center}
\captionof{table}{Coordinate and controller comparison, structural panel.}
\label{tab:coordinate-comparison-structural}
\scriptsize
\begin{tabularx}{\linewidth}{lXXXXXX}
\toprule
Method & physical interpretation & transformation type & spatial or abstract
& linear/nonlinear & fixed/state dependent & global inverse / singularities\\
\midrule
\(dq\) CVC & rotor flux/saliency axes & rotation + current loops & spatial &
linear for given angle & rotor-fixed & global rotation; angle unavailable at standstill without sensing\\
\(ft\) DFVC & flux magnitude and perpendicular torque current & flux-angle rotation & spatial &
linear rotation, nonlinear angle & state-dependent & undefined at \(\psi_s=0\); observer angle error\\
Active flux & equivalent torque-producing flux & flux subtraction/model map & spatial variable + task meaning &
linear only for constant inductances & fixed or map-dependent & weak at \(\psi_a=0\); parameter/map ambiguity\\
Torque/flux IOFL & direct output rates & input transformation & abstract outputs/virtual inputs &
nonlinear & state-dependent & local if \(\det(J_hB)\ne0\); voltage blow-up near rank loss\\
\((i_T,\alpha,z_0)\) & torque amplitude and two allocations & state chart & abstract &
nonlinear & fixed for linear EESM & branchwise; singular at zero torque/factors\\
\((M,P,U^2)\) & performance levels themselves & state chart & abstract &
nonlinear & state-dependent & local only; gradient dependence and multiple branches\\
local gradient frame & torque, voltage, remaining direction & differential basis & abstract &
nonlinear & moving & generally not integrable; critical points/eigen crossings\\
AJD frame & least average quadratic coupling & orthogonal offline rotation & abstract &
linear & fixed & global linear inverse; nonzero residual coupling\\
\bottomrule
\end{tabularx}
\end{center}
\end{landscape}

\begin{landscape}
\begin{center}
\captionof{table}{Coordinate and controller comparison, implementation panel.
Costs refer to the architecture, not only the coordinate transform.}
\label{tab:coordinate-comparison-implementation}
\scriptsize
\begin{tabularx}{\linewidth}{lXXXXXXX}
\toprule
Method & observer & loops & reference maps & FPGA / MCU cost & saturation
handling & field weakening & EESM suitability\\
\midrule
\(dq\) CVC & rotor angle, currents & \(2+1\) PI & MTPA/FW/limits & low / low &
calibrated maps mature & indirect but mature & good baseline, hides redundancy\\
\(ft\) DFVC & stator flux and angle & flux + torque/current + field & fewer FW maps &
pipeline-friendly / medium & nonlinear flux observer & excellent direct flux limit & demonstrated; field allocation remains\\
Active flux & active-flux estimate or rotor angle & usually current/torque loops & reduced torque model &
low--medium / low--medium & map-dependent \(\psi_a\) & helpful, not by itself a limit law & excellent torque unification\\
Torque/flux IOFL & task states and drift & task loops & few steady maps &
matrix pipeline / high & accurate local model needed & direct if flux output & strong if decoupling rank persists\\
\((i_T,\alpha,z_0)\) & currents, parameters & three hierarchical loops & small allocation policy &
exp/log or LUT / medium--high & local map/Jacobian & needs voltage-aware \(\alpha,z_0\) & natural static/dynamic allocation\\
\((M,P,U^2)\) & all performance models & three task loops & potentially none locally &
solve-heavy / high & nonlinear Jacobian & direct voltage coordinate & diagnostic more than global controller\\
local gradient frame & flux Jacobian/gradients & hierarchical & local model tables &
parallel products / high & best local adaptation if map accurate & direct tangent direction & powerful but complex\\
AJD frame & currents, rotor angle & transformed current loops & conventional limits &
very low / low & only as calibration family captures it & residual coupling remains & attractive fixed compromise\\
\bottomrule
\end{tabularx}
\end{center}
\end{landscape}

\begin{center}
\captionof{table}{What the principal representations expose across synchronous-machine types.}
\label{tab:machine-representation-map}
\small
\begin{tabularx}{\linewidth}{lXXXX}
\toprule
Representation & SPM & IPMSM & SynRM & EESM\\
\midrule
\(dq\) currents & nearly direct torque/flux split & saliency remains coupled & reluctance product hidden & field/stator allocation hidden\\
\(ft\) frame & flux limit and torque current & flux limit and total torque current & flux/torque action & flux/torque action; field allocation remains\\
Active flux & fixed PM flux & PM plus reluctance flux & reluctance-producing flux & field plus reluctance flux\\
Task/null chart & one secondary direction & one secondary direction & one secondary direction & two torque-null tangent directions\\
Performance chart & mainly diagnostic & local loss/voltage chart & local loss/voltage chart & richest three-task chart, but most singularity-prone\\
\bottomrule
\end{tabularx}
\end{center}

\subsection{Do we still need \texorpdfstring{\(dq\)}{dq}?}

As a spatial interface, usually yes: it removes sinusoidal time variation and
fixes rotor anisotropy. As the only internal plant representation, no. As
optimization variables, often no. As controlled outputs, frequently no. For
EESM, no single chart is best across startup, loss optimization, field
weakening, saturation, and the capability boundary. The robust architecture
is layered: regular \(dq\) or flux states at the plant boundary, a
machine-specific model block, task variables for optimization and priorities,
and guarded virtual inputs only where their decoupling matrix is healthy.

\subsection{Explicit answers to the fifteen evaluation questions}

\begin{enumerate}
 \item \textbf{Is \(dq\) the best internal control coordinate?} For an SPM it
 is nearly a torque/flux task chart already. For an IPMSM it remains an
 excellent plant frame but not an ideal torque chart. For an EESM it hides two
 allocation freedoms and is best retained as an inner spatial interface.
 \item \textbf{When is \(ft\) superior?} When stator-flux magnitude directly
 expresses the active voltage limit, especially in field weakening, and a
 sufficiently accurate flux observer is available.
 \item \textbf{When is active flux clearer?} With saliency or controllable
 excitation, because it absorbs reluctance, PM, and field contributions into
 the one flux factor that multiplies torque current.
 \item \textbf{Can an EESM torque current be defined?} Yes: \(i_t\) is exact
 in the stator-flux frame, while \(i_T\) is a nonspatial torque-amplitude
 coordinate in the loss-whitened EESM task chart.
 \item \textbf{Which two complements are most useful?} Below base speed,
 allocation imbalance \(\alpha\) and the loss-orthogonal null coordinate
 \(z_0\). In field weakening, replace or deform one complement toward stator
 flux or voltage utilization.
 \item \textbf{Fixed or moving frame?} Fixed \(dq\), active-flux, or AJD
 frames are preferable while their residual coupling is acceptable. Saturated
 and constraint-changing regions may justify a guarded moving local frame.
 \item \textbf{Where are the singularities?} At zero flux for \(ft\), zero
 active flux for torque division, zero torque/factors for hyperbolic charts,
 dependent performance gradients, decoupling-matrix rank loss, and local
 eigenvalue crossings.
 \item \textbf{Most mathematically elegant?} \((M,P,U^2)\) straightens three
 level-set families, but only locally. Hyperbolic coordinates are the cleanest
 global-on-a-branch expression of the linear EESM torque product.
 \item \textbf{Most physically transparent?} \(ft\) for flux/torque action and
 \((M,\psi_a,\sigma)\) for EESM torque production plus internal sharing.
 \item \textbf{Cheapest on MCU/FPGA?} The measured-angle \(dq\) CVC baseline
 remains cheapest in the stated count model. A fixed AJD frame adds little;
 state-dependent task/IOFL frames cost more unless they remove substantial
 maps or calibration.
 \item \textbf{Most robust under saturation?} In principle a local
 gradient/differential-flux frame; in present practice, calibrated \(dq\) or
 DFVC with a validated nonlinear flux map is safer than an unguarded inverse.
 \item \textbf{Largest calibration reduction?} DFVC removes much of the
 field-weakening reference map; active flux unifies torque calculation.
 Neither removes the saturated magnetic map.
 \item \textbf{Natural architecture for the static optimum?} The
 \((i_T,\alpha,z_0)\) hierarchy: set torque amplitude, drive allocation toward
 the loss/voltage optimum, then use the remaining null freedom.
 \item \textbf{Natural architecture for A-to-B motion?} The same task chart is
 the best warm start because it separates rapid torque acquisition from slow
 allocation, while the final constrained optimization remains in nonsingular
 physical coordinates.
 \item \textbf{A unified machine task coordinate system?} Evidence supports a
 unified \emph{interface}---torque, a flux resource, and priority/limit
 information---but not one globally invertible coordinate formula for SPM,
 IPMSM, SynRM, and EESM. Machine-specific inner charts are physically
 unavoidable.
\end{enumerate}

\begin{quote}\itshape
Coordinates are tools, not physical laws. A useful transformation must expose
meaning, optimality, feasibility, dynamics, observability, controllability, or
implementation---and must state where the displaced complexity reappears.
\end{quote}
